\documentclass[11pt]{article}

\usepackage{amsmath,amssymb}
\usepackage{xurl}
\usepackage[hidelinks]{hyperref}
\setlength{\emergencystretch}{2em}

% Shared commands for notes in docs/paper-gaps/.
%
% Two halves.  The link commands below are project identity: OWNER/REPO is the
% placeholder that scripts/bootstrap_project.py replaces with this project's
% GitHub slug and Pages site.  Everything after them is NOTATION, and notation
% belongs to the source paper: the definitions here are an example set, and a
% project replaces them with the macros of its own paper's style file.  The one
% rule that never changes: a command defined here must mean exactly what the
% source paper's macro means, or a note silently says something else.

\newcommand{\Lean}{\textsc{Lean}}
\newcommand{\leanid}[1]{\textsf{#1}}
\newcommand{\ghissue}[1]{\href{https://github.com/OWNER/REPO/issues/#1}{GitHub issue \##1}}
\newcommand{\ghpr}[1]{\href{https://github.com/OWNER/REPO/pull/#1}{PR \##1}}
% Local issue tracker (issues/NNNN-*.md in this repository); no remote URL.
\newcommand{\localissue}[1]{issue \##1 (\path{issues/})}
\newcommand{\doclink}[1]{\href{https://OWNER.github.io/REPO/docs/#1.html}{\path{#1}}}

% --- Notation: replace with the source paper's own macros --------------------
% \F, \N, \C, \R, \tr, \ot, \eps, \E, \bx, \bu, \bv, \by

\newcommand{\F}{\mathbb{F}}
\newcommand{\N}{\mathbb{N}}
\newcommand{\C}{\mathbb{C}}
\newcommand{\R}{\mathbb{R}}
\newcommand{\tr}{\mathrm{tr}}
\newcommand{\eps}{\epsilon}
\newcommand{\ot}{\otimes}
\newcommand{\E}{\mathop{\bf E\/}}

% bold random variables
\newcommand{\bu}{\boldsymbol{u}}
\newcommand{\bv}{\boldsymbol{v}}
\newcommand{\bx}{{\boldsymbol{x}}}
\newcommand{\by}{\boldsymbol{y}}

% T1 so literal < and > in \texttt placeholders typeset as themselves
\usepackage[T1]{fontenc}

% bra-ket
\usepackage{braket}

% --- Example structured spaces ---
\newcommand{\polyfunc}[3]{\mathcal{P}(#1, #2, #3)}
\newcommand{\polymeas}[3]{\mathrm{PolyMeas}(#1, #2, #3)}
\newcommand{\polysub}[3]{\mathrm{PolySub}(#1, #2, #3)}

% Approximation relations (paper uses \approx_\eps and \simeq_\zeta)
\newcommand{\approxeps}{\approx_{\varepsilon}}
\newcommand{\simgeq}{\simeq_{\zeta}}

\renewcommand{\ghissue}[1]{%
  \href{https://github.com/Dengnifer/MIPStarRE-A/issues/#1}{GitHub issue \##1}}

\title{Quantum Linearity and a Direct Combined-Point Construction}
\author{}
\date{2026-09-04}

\begin{document}
\maketitle

\section{At a glance}

\begin{itemize}
  \item \textbf{Difficulty.} Medium.  The source quotes a fiberwise consequence
  of a state-dependent quantum linearity theorem, but suppresses an average,
  inserts an unnecessary bound on the error, appends an identity at the origin,
  and assumes that its ancillary space can be absorbed into zero-state padding
  of the strategy. The factor of two lost between measurement distance and
  operator distance is not one of these: it is a defect of the provider's own printed statement, treated
  in a separate note.
  \item \textbf{Estimated weight.} The corrected single-family theorem and its
  two distance formulations are proved. A fixed optional Naimark ancilla now
  serves all families of the same index length. Absorption of this ancilla
  into the strategy's original zero-state padding is not proved. Three
  auxiliary obstruction lemmas show that the required reserve cannot be
  inferred for every fixed projective setting.
  \item \textbf{Mathlib/project split.} Finite direct sums, extensions of
  operators, and squared norms are standard linear algebra.  The proved
  Fourier--Naimark stability theorem and the two state-dependent distance
  conventions are project-level results.
  \item \textbf{Key mathematical inputs.} Theorem~10 of
  \cite{Natarajan2017Linearity}, the operator BLR identity, finite averaging,
  and Naimark dilation on a fixed index space. The remaining padding
  assumption is discussed in \ghissue{19}.
  \item \textbf{Combined points.} A separate, completed argument proves the
  combined-point lemma on the original expanded spaces.  It uses a
  field-valued sandwich POVM, finite-character Parseval, an exact overlap
  identity, and projective rounding with a dimension-independent bound
  $K\varepsilon^{1/8}$.  This replaces the binary-linearity argument; it does
  not establish the source's padding assumption. The implementation in
  \ghissue{388} comprises roughly 2,400 lines, predominantly project-local
  measurement arguments using Mathlib finite sums, inner products, and real
  powers.
\end{itemize}

\section{Key theorem forms}

Let $H$ be a finite-dimensional Hilbert space, let $\rho$ be an arbitrary
density operator on $H$, and let $u\mapsto O^u$ be binary observables on $H$,
indexed by $\mathbb F_2^t$.  The cited theorem assumes the averaged
correlation estimate
\[
  \mathop{\mathbb E}_{u,v}
  \operatorname{Re}\operatorname{Tr}\!\left(
    O^u O^v O^{u+v}\rho\right) \geq 1-\delta.
\]
Its Fourier--Naimark argument produces an auxiliary space $H'$, a unit vector
$\mathrm{anc}\in H'$, and binary observables $L^u$ on $H\otimes H'$ satisfying,
with
$\rho'=\rho\otimes|\mathrm{anc}\rangle\langle\mathrm{anc}|$,
\[
  L^uL^v=L^{u+v}
  \quad\text{for every }u,v.
\]
Moreover,
\[
  \mathop{\mathbb E}_{u}
  \operatorname{Tr}\!\left((L^u-O^u\otimes I_{H'})^\dagger
    (L^u-O^u\otimes I_{H'})\,\rho'\right)\leq2\delta.
\]
Equivalently, if $d^{\mathrm{bin}}_{\rho'}$ denotes the distance between the
binary projective measurements associated with the two observables, then
\[
  \mathop{\mathbb E}_{u}
  d^{\mathrm{bin}}_{\rho'}(L^u,O^u\otimes I_{H'})^2\leq\delta.
\]
The provider prints $\delta$ for the first, operator-distance bound.  Its
calculation instead proves the second, measurement-distance bound, which is
one half of the squared operator distance.  The factor-of-two discrepancy and
a counterexample to the printed operator bound are given in
\path{docs/paper-gaps/qpbt_linearity-distance-normalization.tex}, tracked in
\ghissue{124}.
In particular, $L^0=I$: exact linearity gives $(L^0)^2=L^0$, while an
observable is invertible.

The chapter needs a parameterized version.  For each pair $(x,z)$ it applies
the theorem to $(r,s)\mapsto O^{x,z,r,s}$.  If $\delta_{x,z}$ is the mean
squared multiplicative defect on that fiber, the corresponding correlation
error is $\delta_{x,z}/2$.  The corrected theorem therefore gives mean squared
operator distance at most $\delta_{x,z}$, equivalently mean squared binary-
measurement distance at most $\delta_{x,z}/2$.  Since the theorem accepts the
actual nonnegative error on each fiber, no assumption that a fiber error is at
most one, and no good--bad fiber split, is needed. The Fourier--Naimark
construction uses the same optional ancillary Hilbert space and unit vector
for every fiber of fixed index length; it does not identify this extension
with the strategy's original expanded local spaces.

For the combined-point conclusion, take admissible QPBT parameters
$q=2^t,m,M$ and a projective strategy on $H_A\otimes H_B$ with unit state
$\psi$ and winning probability at least $1-\varepsilon$.  Its expanded state is
\[
 \Psi=\psi_{AB}\otimes\Phi_{A'A''}\otimes\Phi_{B'B''},\qquad
 \Phi=q^{-M/2}\sum_{v\in\mathbb F_q^M}|v,v\rangle.
\]
Write $X_a^x$ and $Z_b^z$ for the field-valued expanded point effects, and
$H_i^+=H_i\otimes(\mathbb C^q)^{\otimes M}$ for player $i$'s existing expanded
local space.  For families on the six-register space, put
\[
 d^2(A,B)=\mathop{\mathbb E}_{x,z}\sum_{a,b\in\mathbb F_q}
   \|(A^{x,z}_{a,b}-B^{x,z}_{a,b})\Psi\|^2.
\]
Lemma~\path{lem:qld-4-10}, at lines~689--709 of
\path{references/qpbt-paper/14_analysis_of_the_pauli_basis_test.tex}, asserts
projective measurements $Q_i^{x,z}$ on $H_i^+$ with outcomes in
$\mathbb F_q^2$.  The opposite placement pairs are $(AA',BA'')$ and
$(BB',AB'')$, in either direction.  For each such pair $(p_1,p_2)$ its three
conclusions are
\[
 d^2(Q_{p_1},Q_{p_2}),\quad
 d^2(Q_{p_1},(XZ)_{p_2}),\quad
 d^2(Q_{p_1},(ZX)_{p_2})\ \leq\ \delta_Q(\varepsilon),
\]
where $\delta_Q$ is polynomially small in the paper's convention.  Here
$(XZ)^{x,z}_{a,b}=X_a^xZ_b^z$ and $(ZX)^{x,z}_{a,b}=Z_b^zX_a^x$;
subscripts denote register placement, using the measurement for that player.
The direct construction below gives $\delta_Q(\varepsilon)=K\varepsilon^{1/8}$
for a universal $K$, without another ancillary extension.

\section{Comparison with the quoted statement}

Lines~711--725 of
\path{references/qpbt-paper/14_analysis_of_the_pauli_basis_test.tex} quote the
linearity theorem.  The displayed hypothesis correctly averages over both
group variables.  The displayed conclusion, however, says that the corrected
observable is close for every $u$, whereas Theorem~10 of the cited work gives
closeness only on average over $u$ \cite{Natarajan2017Linearity}.  The
quotation also inserts $\delta\leq1$, which is absent from the provider theorem
and would unnecessarily obstruct applying it with a fiber's actual defect, and
it appends the assertion at the identity, which follows from exact linearity as
above.  These three are the differences between the quotation and the cited
statement.

The factor of two is not one of them.  The provider's equation~(3) proves the
bound $\delta$ for binary-measurement distance; under the operator-distance
convention displayed above this is the bound $2\delta$, and it is the
provider's own Theorem~10 that prints $\delta$ for the operator distance.  The
quotation states its conclusion in the relation $\approx_\delta$, which asserts
only an $O(\delta)$ bound and therefore absorbs the constant factor instead of
reproducing the erroneous one.  The normalization defect and its counterexample
are the provider's and are treated in
\path{docs/paper-gaps/qpbt_linearity-distance-normalization.tex}.

Lines~825--832 construct the family $O^{x,z,r,s}$ and establish approximate
linearity only after averaging over $x,z,r,r',s,s'$. They then apply the
single-family theorem separately for every $(x,z)$ and state that sufficiently
many zero ancillas avoid a further extension. The Fourier--Naimark construction
actually provides one optional ancillary space, with basis indexed by
$\mathbb F_2^{2t}\sqcup\{\bot\}$ and ancilla vector at $\bot$, for every
fiber $(x,z)$. This uniformity does not imply that the fixed state
$\ket{\hat\psi}$ already contains this ancillary register in its zero state.
To follow the source's linearity route on the original strategy spaces one
would still have to justify that padding or prove a state-preserving
absorption theorem. Simply compressing the Naimark projection-valued
measurement along the ancillary vector gives the Fourier-square POVM, which
need not be projective. The completed field-valued proof of the combined-point
lemma avoids this step entirely.

\section{Fixed spaces and reserved subspaces}

The source permits a choice of padded strategy at lines~160--172, recalls
that convention at lines~367--372, and uses it at line~832. This is an
existential choice made before the subsequent measurement constructions.
It is not the assertion that every projective strategy already contains an
unused pure register. The Lean type \path{ProjectiveSetting} retains the
strategy, projectivity, and winning inequality, without specifying a reserved
subspace. The theorem \path{exists_projective_padded_strategy} supplies a
particular zero-padded strategy, but its public conclusion records neither
an unused register nor an invariant smaller space for later observables.

Two elementary obstructions prevent inferring this information from the
dimensions alone. Put $K_t=\mathbb C^{\mathbb F_2^t\sqcup\{\bot\}}$.
For nonzero finite-dimensional $H$,
\[
 \dim(H\otimes K_t)=(2^t+1)\dim H>\dim H,
\]
so the full extension returned by linearity on $H$ cannot embed isometrically
back into $H$. Moreover, a positive definite local state cannot become
$\sigma\otimes|\bot\rangle\langle\bot|$ under an invertible coordinate
change: positive definiteness is preserved, whereas that tensor product
annihilates every Boolean-word ancillary coordinate. Thus even an adequate
ambient dimension does not imply the required state identity.

There is also a counterexample within the actual projective-setting domain.
For any admissible parameters, take the honest projective strategy on its
EPR state of local dimension $D$, and let $\varepsilon$ be its rejection
probability. The winning premise holds and $\varepsilon\geq0$. Its expanded
state has reduced density $I/(Dq^M)$ on $AA'$, where $M=2^m$. Consequently,
\[
 R_{AA'}\Psi=\Psi\quad\Longrightarrow\quad R=I.
\]
In particular, no proper projection onto a zero ancillary sector fixes
this state, even after changing its local orthonormal basis. The Lean proof
compares the coefficients of the three EPR factors directly, without
introducing a partial-trace definition. It asserts a nonnegative rejection
probability, not a separate value-one theorem for this witness.
These facts are verified by
\path{not_nonempty_linearity_ancilla_isometry},
\path{posDef_ne_linearity_ancilla_conjugate}, and
\path{exists_projectiveSetting_no_proper_alice_support} in
\path{MIPStarRE/QPBT/Combining/Linearity/Padding.lean}. They refute automatic
reservation for arbitrary fixed settings, not the source's advance choice
of padding and not its combined-point conclusion.

A sufficient construction must keep that advance choice explicit. For
$q=2^t$ the common ancilla for the pairs $(r,s)$ has dimension $q^2+1$;
$2t+1$ zero qubits suffice to contain it. This number is independent of
$(x,z)$, the observables, the state, and the error. It does not change the
Pauli-register parameter $M=2^m$. One must construct a projective strategy
with an invariant smaller local space $H_0$, and an isometry
$J:H_0\otimes K_{2t}\longrightarrow H^+$ into its already chosen expanded
space, with the original state embedded by $v\mapsto J(v\otimes\bot)$.
The original measurements must intertwine with this embedding. Every
subsequent observable family must be constructed on $H_0$ and transported
using this same embedding; this is stronger than assuming a transport
certificate for a family chosen on all of $H^+$.

There is a concrete candidate in the existing setup construction. If $N$ is
the size of its answer alphabet, the active Naimark summand has local index
$I\times\operatorname{Option}(\mathrm{Answer})$, of dimension
$|I|(N+1)$, inside a Boolean-padded space of dimension $|I|2^{N+1}$.
The measurements preserve this summand and the state lies in it; the
complement carries a fixed-outcome measurement. For the Pauli answer alphabet,
the Pauli-outcome constructor gives $N\geq q^{2^m}\geq q^2\geq4$, so
\[
 (N+1)(q^2+1)\leq(N+1)^2\leq2^{N+1}.
\]
Hence the existing Boolean cube has enough dimension to embed the active
summand times the common ancilla. The required embedding must agree on the
ground slice with the existing one-hot embedding; the inequality alone
does not construct that agreement or the measurement transport.

After such an embedding is constructed, a rounded observable $L^u$ extends
to $H^+$ as $JL^uJ^\dagger+(I-JJ^\dagger)$. Orthogonality of the two
summands preserves exact linearity and binary observability. Intertwining
of the original operators on the ground slice preserves each squared
state-dependent error, and hence both the Boolean-index and $(x,z)$ averages.
For the corresponding projective measurement the complementary projection
is assigned to the zero outcome. This is a proposed sufficient construction,
not a proved absorption theorem. Its state identity, all-family transport,
and use in the binary construction remain open under \ghissue{697}.

\section{Direct construction of the combined points}
\label{sec:direct-combined-points}

The source's three-step plan at line~711 and proof at lines~731--881 first
combine binary refinements, prove approximate linearity in their scalar
indices, and apply quantum linearity.  The following replacement retains the
statement of \path{lem:qld-4-10}, but forms and rounds a sandwich with
field-valued outcomes.  The required commutation estimate is obtained without
a factor depending on the field size.

\paragraph{Parseval transfer.}
Fix a point pair and a placement, suppressing these indices.  With
$\chi(v)=(-1)^{\operatorname{tr}(v)}$, the Fourier formulas at paper
lines~384--418 give
\[
 X_a=q^{-1}\sum_r\chi(ar)\widehat X^r,\qquad
 Z_b=q^{-1}\sum_s\chi(bs)\widehat Z^s.
\]
Write $[A,B]=AB-BA$.  Character orthogonality applied to the vectors
$[\widehat X^r,\widehat Z^s]\Psi$ yields the exact identity
\[
 \sum_{a,b}\|[X_a,Z_b]\Psi\|^2
 =q^{-2}\sum_{r,s}\|[\widehat X^r,\widehat Z^s]\Psi\|^2.
\]
Indeed, the commutator Fourier expansion has coefficient $q^{-2}$; squaring
contributes $q^{-4}$, and orthogonality contributes $q^2$.  The remaining
$q^{-2}$ is precisely uniform averaging.  This is finite-dimensional Parseval
for vector-valued sums, not an estimate of each outcome separately.  The
observable commutation bound underlying \path{lem:qld-comm-cons}, item~2
(paper lines~466--505 and its proof), therefore gives
$\mathbb E_{x,z}\sum_{a,b}\|[X_a,Z_b]\Psi\|^2\leq C_{\rm com}\sqrt\varepsilon$
on each placement.  Equivalently, the binary-refinement commutators are
$(-1)^{c+d}[\widehat X^r,\widehat Z^s]/4$, so passing from their summed
squared norms to the observable bound costs only a factor of four.

\paragraph{The sandwich and its overlap.}
For each player and point pair set $R_{a,b}=Z_bX_aZ_b$.
Projectivity gives $R_{a,b}=(X_aZ_b)^\dagger(X_aZ_b)\geq0$ and
$\sum_{a,b}R_{a,b}=I$.  Moreover,
$R_{a,b}-Z_bX_a=Z_b[X_a,Z_b]$, so contraction by $Z_b$ shows
$d^2(R,ZX)\leq C_{\rm com}\sqrt\varepsilon$.

Now fix opposite placements, denoted by subscripts $1,2$, and a point pair.
Every operator on the first placement commutes with every operator on the
second.  Put
\[
 D^X_a=X_{1,a}-X_{2,a},\quad D^Z_b=Z_{1,b}-Z_{2,b},\quad
 W_b=Z_{1,b}Z_{2,b},\quad K^i_{a,b}=[X_{i,a},Z_{i,b}].
\]
For complete projective measurements $P,Q$ on opposite placements and any
vector $w$, expanding the squared differences gives
\[
 \sum_c\langle w,P_cQ_cw\rangle
 =\|w\|^2-\tfrac12\sum_c\|(P_c-Q_c)w\|^2.
\]
Since $R_{1,a,b}R_{2,a,b}=W_b^\dagger X_{1,a}X_{2,a}W_b$, apply this identity
first to $X_1,X_2$ in $W_b\Psi$, and then to $Z_1,Z_2$ in $\Psi$.  The result is
\[
 1-\sum_{a,b}\langle\Psi,R_{1,a,b}R_{2,a,b}\Psi\rangle
 =\tfrac12\sum_b\|D^Z_b\Psi\|^2
   +\tfrac12\sum_{a,b}\|D^X_aW_b\Psi\|^2.
\]
The left side is the off-diagonal consistency mass because both POVMs are
complete.  This exact identity replaces the source's Cauchy--Schwarz chain
\path{eq:qld-rw-self-cons-1} through \path{eq:qld-rw-self-cons-4}
(paper lines~743--778).

To bound its second term, use the algebraic decomposition
\[
 D^X_aW_b=W_bD^X_a+Z_{2,b}K^1_{a,b}-Z_{1,b}K^2_{a,b}.
\]
The squared norm of a sum of three vectors is at most three times the sum
of their squared norms.  Projection contraction bounds the last two terms,
and $\sum_b\|W_bv\|^2\leq\sum_b\|Z_{2,b}v\|^2=\|v\|^2$ bounds the first
after summing over $b$.  Thus the off-diagonal mass is at most
\[
 \begin{aligned}
 c(x,z):={}&\tfrac12\sum_b\|D^Z_b\Psi\|^2
       +\tfrac32\sum_a\|D^X_a\Psi\|^2\\
   &+\tfrac32\left(\sum_{a,b}\|K^1_{a,b}\Psi\|^2
       +\sum_{a,b}\|K^2_{a,b}\Psi\|^2\right).
 \end{aligned}
\]
Item~1 of \path{lem:qld-comm-cons} bounds each averaged point-consistency
sum by $C_{\rm sc}\varepsilon$.  Together with Parseval this gives
\[
 \mathbb E_{x,z}c(x,z)\leq
 2C_{\rm sc}\varepsilon+3C_{\rm com}\sqrt\varepsilon
 \leq C_0(\varepsilon+\sqrt\varepsilon),
 \qquad C_0=2C_{\rm sc}+3C_{\rm com}.
\]
All constants can be taken at least one and independent of $q,m,M$ and the
strategy dimensions.

\paragraph{Rounding on the original spaces.}
The orthonormalization lemma \path{lem:ortho}, also used at paper
lines~782--786, says that a local POVM consistent with a POVM on the other
tensor factor can be replaced by a projective measurement on the same local
space.  Its proved explicit form gives summed squared distance at most
$220c^{1/4}$ whenever the consistency defect is at most $c\geq0$; no upper
bound on $c$ is required.  Apply it separately to each point pair, with
bipartitions $AA'\mid BA''(B'B'')$ for Alice and
$BA''\mid AA'(B'B'')$ for Bob.  All choices lie in the already fixed spaces
$H_A^+$ and $H_B^+$, so no compatibility of newly returned spaces is needed.
Concavity of the fourth root gives
\[
 d^2(Q,R)\leq220\,\mathbb E_{x,z}c(x,z)^{1/4}
 \leq220\,[C_0(\varepsilon+\sqrt\varepsilon)]^{1/4}.
\]
The ancillary permutation $(A'\ B'')(A''\ B')$ fixes $A,B$ and fixes
$\Psi$: it exchanges the two identical EPR factors and reverses both pairs.
It transports these same-side distances from $AA'$ to $AB''$ and from $BA''$
to $BB'$.  This independently proved involution requires no symmetric
strategy.  It is not the four-cycle $(A'\ B''\ A''\ B')$ induced by
$\sigma\theta$ in blueprint \path{lem:symmetric-equivalents-transfer}.

\paragraph{The final error.}
Transporting one factor at a time across opposite placements gives
$d^2((ZX)_{p_1},(XZ)_{p_2})\leq C_3\varepsilon$ with $C_3=4C_{\rm sc}$.
The contraction used in this transport is
$\sum_a\|M_av\|^2\leq\|v\|^2$, valid for every POVM since
$\sum_a M_a^\dagger M_a\leq\sum_a M_a=I$.
Let $C_1=C_2=C_{\rm com}$ bound respectively the sandwich-to-product and
commutation errors, and take the convenient larger rounding bound
$\eta=440[C_0(\varepsilon+\sqrt\varepsilon)]^{1/4}$.
Repeatedly using $d^2(A,C)\leq2d^2(A,B)+2d^2(B,C)$ along the sequence
\[
 Q_{p_1},\ R_{p_1},\ (ZX)_{p_1},\ (XZ)_{p_2},\ (ZX)_{p_2},\ R_{p_2},\ Q_{p_2}
\]
bounds all three required distances by
\[
 12\eta+20C_1\sqrt\varepsilon+16C_3\varepsilon+4C_2\sqrt\varepsilon.
\]
For $0\leq\varepsilon\leq1$, both $\varepsilon$ and $\sqrt\varepsilon$ are
at most $\varepsilon^{1/8}$, while
$[C_0(\varepsilon+\sqrt\varepsilon)]^{1/4}
 \leq2C_0\varepsilon^{1/8}$.
Hence one may take $K=12(880C_0)+20C_1+16C_3+4C_2$.
For $\varepsilon>1$, each relevant family has
$\sum_{a,b}A_{a,b}^\dagger A_{a,b}\leq I$ (equality for the ordered products
of projective measurements), so every distance is at most $4\leq
K\varepsilon^{1/8}$.  No step incurs a factor from the number of outcomes,
point pairs, or Hilbert-space dimension.

\section{Blueprint and formalization status}

Theorem~\path{thm:linearity} and Remark~\path{rem:linearity-import} in
\path{blueprint/src/chapter/ch15_qpbt_combining.tex} record respectively the
averaged provider theorem with the corrected distance normalization and the
distinction between a uniform optional ancilla and its absorption into a
fixed strategy. Lemma~\path{lem:linearity-common-ancilla} records the first
fact as a formalization-only consequence of the Fourier--Naimark construction.

The declaration \path{exists_exactly_linear_observables} in
\path{MIPStarRE/QPBT/Combining/Linearity.lean} assumes a nonnegative defect, a
positive semidefinite trace-one operator, binary observables, and the averaged
three-observable correlation.  It returns one finite ancillary space and
normalized ancillary vector for that family, exactly linear binary observables
on the extension, and the averaged squared operator-distance bound $2\delta$.
The companion declaration
\path{exists_exactly_linear_observables_binaryObservableDistSq} gives the
equivalent squared binary-measurement bound $\delta$.  Both declarations are
proved.  They neither assume $\delta\leq1$ nor assert pointwise closeness, and
the first declaration records the correction to the provider's printed
operator bound.  Both are linked from Theorem~\path{thm:linearity}.  The
normalization analysis was carried out in \ghissue{124} and the statement
repair in \ghissue{129}; both are closed. The additional declaration
\path{exists_exactly_linear_observables_commonAncilla} selects the ancillary
space with basis indexed by $\mathbb F_2^t\sqcup\{\bot\}$ and its canonical
unit vector at $\bot$ before quantifying over the finite Hilbert space, the
density operator, the error, and the observable family. It returns exact
linearity for all pairs and mean
squared operator distance at most $2\delta$ for each eligible family. It
does not supply zero-state padding on the existing expanded player spaces.
Independently, \path{exists_combinedPointsWitness} in
\path{MIPStarRE/QPBT/Combining/Points.lean} proves the combined-point statement
of blueprint \path{lem:qld-4-10} on the original expanded spaces by
Section~\ref{sec:direct-combined-points}.  The modules in
\path{MIPStarRE/QPBT/Combining/Points/} formalize respectively placement
invariance, Parseval commutation transfer, the sandwich POVM, its overlap
identity and consistency bound, projective rounding, and the final distance
estimates.  This is a replacement proof of the same statement, not a
formalization of the source's binary-linearity proof. It neither invokes nor
discharges the zero-state-padding assumption of \path{rem:linearity-import}.
The affine coarse-graining of these constructed joint point measurements,
including self-consistency and both ordered point-product comparisons with
the same universal error, is the declaration \path{exists_extendedQ} for
source Lemma~\path{lem:qld-4-12}. It does not establish the subsequent
extended-line lemma \path{lem:qld-4-13}.
The complete QPBT audit status is recorded in
\path{docs/paper-gaps/qpbt-gap-register.md}, whose linearity row records the
uniform-ancilla result together with the residual zero-state-padding
assumption.

Issue~\#19 introduced this discrepancy record for the extraction skeleton;
\ghissue{47} owns the Stage~4.3 proof work consumed downstream.  The chapter-16
extraction declarations consume the concrete \path{GlobalPairWitness}; they
do not assume a new fiberwise linearity hypothesis or a common ancillary space.

\section{Conclusion}

The quotation changes averaged closeness to pointwise closeness, inserts an
unnecessary upper bound on $\delta$, and appends an assertion at the identity
that is derivable from exact linearity.  It does not reproduce the provider's
factor-of-two normalization defect, since it states its conclusion in a
relation that absorbs constant factors; that defect is the provider's own and
is treated separately. The single-family theorem is proved with mean squared
operator error at most $2\delta$, equivalently mean squared binary-measurement
error at most $\delta$. Its optional ancillary space and unit vector can be
chosen uniformly over the fibers indexed by $(x,z)$, but their absorption
into the fixed strategy's zero-state padding is unproved. The combined-point
conclusion is proved independently on the original expanded spaces, with
universal error $K\varepsilon^{1/8}$ and no additional hypotheses; its scalar
coarse-graining establishes the extended-point lemma with the same error.
The extended-line lemma remains open. There is no statement discrepancy in
the replacement point construction.

\bibliographystyle{alpha}
\bibliography{references}

\end{document}
